\appendix
\section{Reproducibility Manifest}
\label{sec:appendix_manifest}

Every cited number in this paper traces to a specific entry in the
per-run SHA-256 manifest published alongside the code:

\begin{itemize}
\item \texttt{paper2/runs/<run\_id>/MANIFEST.json} --- the
      definitive index of every artefact produced by the run, with
      file size, sha256, and producer-stage annotation.
\item \texttt{paper2/runs/<run\_id>/metrics/combined\_metrics.json}
      --- the flat dictionary the LaTeX layer consumes via
      \texttt{report.inject.\textbackslash VAR\{\}}.
\item \texttt{paper2/runs/<run\_id>/env/git\_sha.txt} --- the git
      SHA of the producing tree.
\item \texttt{paper2/runs/<run\_id>/env/config\_snapshot.json} ---
      the exact \texttt{ExpConfig} dataclass that produced the run.
\item \texttt{paper2/runs/<run\_id>/env/pip\_freeze.txt} --- the
      Python environment.
\item \texttt{paper2/runs/<run\_id>/env/nvidia\_smi.txt} --- the
      GPU model and driver version.
\end{itemize}

The reproducibility recipe is a single command from a clean clone:

\begin{verbatim}
git clone <repo>
cd kaggle2
python paper2/main_paper2.py --stage all
\end{verbatim}

\section{Per-bug F1 Atlas}
\label{sec:appendix_bug_atlas}

The per-bug $\Delta$F1 numbers in Sec.~\ref{sec:bugs} are sourced
from the multi-seed bug-ablation runner.  The atlas table
(\texttt{paper2/results/bug\_atlas.json}) reports per-seed
$\Delta$F1 plus paired-bootstrap CI for every bug guard.
